\section{Background: tensors, factorization, and three components}
\label{sec:tensor-primer}

\subsection{Scalar, vector, matrix, tensor}

The word ``tensor'' sounds intimidating, but here it means a \emph{multi-dimensional array}. The
table below contains almost all the background needed for the paper's formulas.

\begin{center}
\small
\begin{tabular}{@{}llll@{}}
  \toprule
  Object & Order & Example notation & How to read it \\
  \midrule
  Scalar & 0 & $x=7$ & one number \\
  Vector & 1 & $\mathbf{w}=(w_1,\ldots,w_k)$ & a list over keywords \\
  Matrix & 2 & $\mathbf{W}\in\mathbb{R}^{k\times l}$ & a keyword $\times$ location table \\
  Tensor & 3 & $\mathcal{X}\in\mathbb{R}^{k\times l\times L_c}$ & a stack of tables over time \\
  \bottomrule
\end{tabular}
\end{center}

Indices follow axis order from left to right. For a third-order tensor, $x_{ij}(t)$ is the entry at
keyword $i$, location $j$, time $t$. If $k=2,l=2,L_c=3$, the tensor has $2\cdot2\cdot3=12$ entries.
\emph{Unfolding} (matricization) turns a tensor into a matrix by grouping axes; unfolding along the
keyword axis, for example, produces a matrix that linear algebra can process. We do not need to
unfold by hand, but we should remember that it only \emph{rearranges} numbers and creates no data.

\subsection{Tucker factorization, intuitively}

Suppose there are $k$ keywords and $l$ locations. Instead of storing every tensor value, assume the
data are generated by a few shared \emph{styles of dynamics}. Let $d_k$ be the number of keyword
groups and $d_l$ the number of location groups, with $d_k\ll k$ and $d_l\ll l$. The trend is then
approximated by
\begin{equation}
  \widehat{\mathcal{X}}_d
  =\mathcal{W}^{(\mathrm{core})}
   \times_{\mathrm{key}}\mathbf{W}^{(\mathrm{key})}
   \times_{\mathrm{loc}}\mathbf{W}^{(\mathrm{loc})},
  \label{eq:tucker-intuition}
\end{equation}
where $\mathcal{W}^{(\mathrm{core})}$ is a small $d_k\times d_l\times L_c$ tensor, and
$\mathbf{W}^{(\mathrm{key})}$ and $\mathbf{W}^{(\mathrm{loc})}$ are non-negative coefficient
matrices of size $d_k\times k$ and $d_l\times l$.

The symbols $\times_{\mathrm{key}},\times_{\mathrm{loc}}$ denote \emph{mode products}. Read
Equation~\eqref{eq:tucker-intuition} as follows: take the small core, spread it over keyword groups
with weights $\mathbf{W}^{(\mathrm{key})}$, then over location groups with
$\mathbf{W}^{(\mathrm{loc})}$. A keyword with a large weight in group 2 receives more of group 2's
dynamics. Non-negative weights cannot flip a dynamics' sign, so ``this group is rising'' remains
readable after mixing.

\begin{keypoint}
Tucker factorization is not arbitrary file compression. It yields a small core containing
\emph{shared dynamics} and factors saying \emph{which members belong to which group}. This is the
bridge between forecasting and interpretation.
\end{keypoint}

\subsection{A \texorpdfstring{$2\times2\times3$}{2x2x3} example and the benefit of compression}

Imagine two keywords, two locations, and three weeks. The raw tensor has 12 numbers. If we assume one
keyword group and one location group ($d_k=d_l=1$), the core has $1\cdot1\cdot3=3$ numbers and the
two factors have $1\cdot2+1\cdot2=4$: seven parameters instead of 12. On real data the gap is much
larger. A Device window has $8\cdot50\cdot104=41{,}600$ entries. With $d_k=3,d_l=2$, the core has
$3\cdot2\cdot104=624$ entries and the factors have $3\cdot8+2\cdot50=124$: about 748 parameters,
less than 2\% of the window.

\begin{plainwords}
Compression here is not about saving disk space. It forces the model to explain tens of thousands of
observations with a few dynamics and a table of ``who belongs to which group''. If the model still
fits, that group structure is a useful hypothesis.
\end{plainwords}

\subsection{Trend, seasonal, and outlier in a tiny series}

The three components in D-Tracker add up to approximately the data:
\begin{equation}
  \mathcal{X}^{c}\approx\widehat{\mathcal{X}}_d
  +\widehat{\mathcal{X}}_s+\widehat{\mathcal{X}}_o.
  \label{eq:decompose-primer}
\end{equation}
Table~\ref{tab:three-components} is a numerical example \textbf{created for this guide}. The
``Observed'' row is exactly the sum of the other three; a dash in the outlier row means zero.

\begin{table}[htbp]
  \centering
  \small
  \caption{Illustrative trend + seasonal + outlier decomposition.}
  \label{tab:three-components}
  \begin{tabular}{lrrrrrrrr}
    \toprule
    Component & Week 1 & 2 & 3 & 4 & 5 & 6 & 7 & 8 \\
    \midrule
    Observed  & 10 & 12 & 15 & 14 & 11 & 13 & 17 & 25 \\
    Trend     & 10 & 11 & 12 & 13 & 14 & 15 & 16 & 17 \\
    Seasonal  &  0 &  1 &  3 &  1 & $-3$ & $-2$ &  1 &  3 \\
    Outlier   &  -- & -- & -- & -- & -- & -- & -- &  5 \\
    \bottomrule
  \end{tabular}
\end{table}

The trend changes relatively smoothly. The seasonal part repeats with a period (a short pattern is
shown here; real data use weekly/annual periods). The outlier part is unusual and assumed
\emph{sparse}: only a few cells are non-zero. This separation prevents a one-off spike from being
mistaken for a permanent change in trend.

\subsection{From decomposition to the paper's question}

We now have three pieces: an observed tensor, a way to compress its trend into a core and factors,
and a three-part decomposition. The paper adds two more ideas:
\begin{itemize}
  \item the trend core does not evolve arbitrarily; it follows a \textbf{reaction--diffusion} system; and
  \item the model is updated as a stream arrives, with the latent structure and switch times selected
        automatically.
\end{itemize}
Section~\ref{sec:reaction-diffusion} explains the first; Section~\ref{sec:model} assembles the full model.

\begin{commonmistake}
Three common errors: (1) treating $\mathcal{W}^{(\mathrm{core})}$ as a shrunken copy of the raw data ---
it is a \emph{latent parameter} that is learned; (2) treating $\mathbf{W}$ as probabilities --- they
are non-negative weights and need not sum to one; (3) assuming the decomposition is unique --- tensor
factorizations are generally only identifiable up to transformations, and the paper solves a
particular optimization problem.
\end{commonmistake}

\begin{selfcheck}
\begin{enumerate}
  \item With $d_k=2$, $d_l=3$, and $L_c=10$, how many entries are in the core?
  \item Explain in one sentence why non-negative factors preserve the direction of a dynamics.
  \item In the three-component table, what would the model misunderstand if the outlier at week 8 were removed?
\end{enumerate}
\end{selfcheck}
